\section{K-injective limits and finite cohomological dimension}\label{proofs-9746-10132-k-injective-limits-and-finite-cohomological-dimension}

Private source-order review of original derived.tex 9746–\allowbreak{}10132. Authority:\allowbreak{} Stacks a04446e5\allowbreak{}7ec1fbc2\allowbreak{}52a871af\allowbreak{}cec7752f\allowbreak{}b2807b14. Receiving public edition:\allowbreak{} 95a51593\allowbreak{}aceb247a\allowbreak{}c7a18311\allowbreak{}1678c3a0\allowbreak{}f8b8f4b5. The accompanying source-reading and operation records bind the files. Editorial consequences are additional statements with complete arguments; they are not substituted for the original author\textquotesingle s hypotheses. No novelty is claimed.

Human source:\allowbreak{} \href{https://github.com/stacks/stacks-project/blob/a04446e57ec1fbc252a871afcec7752fb2807b14/derived.tex\#L9746}{The Stacks Project,\allowbreak{} original Derived Categories TeX},\allowbreak{} with its cited Homology results. This review continues the full Hom-complex and shift calculation in PROOFS\_\allowbreak{}9206\_\allowbreak{}9744.md,\allowbreak{} Section 7.

\subsection{1. Why a K-injective complex computes every right derived functor}\label{proofs-9746-10132-why-a-k-injective-complex-computes-every-right-derived-functor}

Let I be K-injective and let s:\allowbreak{}I\allowbreak{}→N be a quasi-isomorphism in K(A)\allowbreak{}. Precomposition gives a bijection Hom\_\allowbreak{}K(N,\allowbreak{}I)\allowbreak{}\allowbreak{}→Hom\_\allowbreak{}K(I,\allowbreak{}I)\allowbreak{}. Thus there exists a unique r:\allowbreak{}N\allowbreak{}→I with r s\allowbreak{}=id\_\allowbreak{}I in K(A)\allowbreak{}. It is also a quasi-isomorphism:\allowbreak{} H(r)\allowbreak{} is inverse to H(s)\allowbreak{}. Therefore id\_\allowbreak{}I is a final object of I/\allowbreak{}Qis,\allowbreak{} including the uniqueness required for finality,\allowbreak{} not only the existence of a retraction.

For an arbitrary functor F on K(A)\allowbreak{},\allowbreak{} the system sending s:\allowbreak{}I\allowbreak{}→N to F(N)\allowbreak{} has value F(I)\allowbreak{} at that final object and is an essentially constant ind-object. Indeed the unique arrows to the final object provide its cocone,\allowbreak{} and the Hom colimit from any object is the Hom to F(I)\allowbreak{}. In the source\textquotesingle s exact-functor setting this proves I computes RF. If K is isomorphic in D(A)\allowbreak{} to I,\allowbreak{} the isomorphism K\allowbreak{}→I is represented by an actual map in K(A)\allowbreak{},\allowbreak{} since maps to a K-injective complex agree in K and D. Its induced cohomology maps are isomorphisms. Hence the computation transfers to K by that quasi-isomorphism.

If every complex maps by a quasi-isomorphism to a K-injective complex,\allowbreak{} the computation criterion gives RF everywhere. A quasi-isomorphism s:\allowbreak{}I\allowbreak{}→J between K-injectives is an isomorphism already in K:\allowbreak{} the unique retraction r satisfies r s\allowbreak{}=id\_\allowbreak{}I,\allowbreak{} and precomposition by s is injective on Hom\_\allowbreak{}K(J,\allowbreak{}J)\allowbreak{},\allowbreak{} so (s r)\allowbreak{}s\allowbreak{}=s\allowbreak{}=id\_\allowbreak{}J s implies s r\allowbreak{}=id\_\allowbreak{}J. Every functor F therefore takes this map to an isomorphism. This proves the two source lemmas without imposing injectivity on individual terms.

\subsection{2. The split inverse system and compatible primitives}\label{proofs-9746-10132-the-split-inverse-system-and-compatible-primitives}

Let …\allowbreak{}→I\_\allowbreak{}2\allowbreak{}→I\_\allowbreak{}1 be the source tower,\allowbreak{} with K-injective stages,\allowbreak{} split surjections in each degree,\allowbreak{} and existing termwise limit I. For an acyclic M let D\_\allowbreak{}n\allowbreak{}=Hom\textasciicircum{}•(M,\allowbreak{}I\_\allowbreak{}n)\allowbreak{},\allowbreak{} using the full differential

D\_\allowbreak{}n\textasciicircum{}r\allowbreak{}=∏\_\allowbreak{}a Hom(M\textasciicircum{}a,\allowbreak{}I\_\allowbreak{}n\textasciicircum{}\{a\allowbreak{}+r\})\allowbreak{},\allowbreak{} (d f)\allowbreak{}\textasciicircum{}a\allowbreak{}=d\_\allowbreak{}\{I\_\allowbreak{}n\}\textasciicircum{}\{a\allowbreak{}+r\}f\textasciicircum{}a-(-1)\allowbreak{}\textasciicircum{}r f\textasciicircum{}\{a\allowbreak{}+1\}d\_\allowbreak{}M\textasciicircum{}a.

Every D\_\allowbreak{}n is acyclic by the all-shift K-injectivity calculation in PART14. The transition maps on D\_\allowbreak{}n\textasciicircum{}r are surjective:\allowbreak{} lift each component through the chosen split epimorphism of I terms. Representable Hom and products preserve limits,\allowbreak{} and the differentials agree after each projection. Consequently Hom\textasciicircum{}•(M,\allowbreak{}I)\allowbreak{}\allowbreak{}=lim\_\allowbreak{}n D\_\allowbreak{}n as actual complexes of abelian groups.

Here is an explicit proof of its degree-zero exactness. A compatible cycle z\allowbreak{}=(z\_\allowbreak{}n)\allowbreak{} admits primitives h\_\allowbreak{}n\allowbreak{}∈D\_\allowbreak{}n\textasciicircum{}\{-1\}. Choose h\_\allowbreak{}1 first. Given a compatible primitive through stage n,\allowbreak{} choose any k\allowbreak{}∈D\_\allowbreak{}\{n\allowbreak{}+1\}\textasciicircum{}\{-1\} with d k\allowbreak{}=z\_\allowbreak{}\{n\allowbreak{}+1\}. Its discrepancy t\_\allowbreak{}n(k)\allowbreak{}-h\_\allowbreak{}n is a cycle in D\_\allowbreak{}n\textasciicircum{}\{-1\},\allowbreak{} where t\_\allowbreak{}n is the transition. Acyclicity gives y\_\allowbreak{}n\allowbreak{}∈D\_\allowbreak{}n\textasciicircum{}\{-2\} with d y\_\allowbreak{}n\allowbreak{}=t\_\allowbreak{}n(k)\allowbreak{}-h\_\allowbreak{}n. Lift y\_\allowbreak{}n to y\allowbreak{}∈D\_\allowbreak{}\{n\allowbreak{}+1\}\textasciicircum{}\{-2\} by surjectivity and put h\_\allowbreak{}\{n\allowbreak{}+1\}\allowbreak{}=k-d y. Then d h\_\allowbreak{}\{n\allowbreak{}+1\}\allowbreak{}=z\_\allowbreak{}\{n\allowbreak{}+1\} and t\_\allowbreak{}n(h\_\allowbreak{}\{n\allowbreak{}+1\})\allowbreak{}\allowbreak{}=h\_\allowbreak{}n. The resulting compatible primitive proves H\^{}0(lim D\_\allowbreak{}n)\allowbreak{}\allowbreak{}=0,\allowbreak{} and hence I is K-injective. Shifting this argument proves exactness in every degree.

This is exactly the use of the degree -2 system in the source\textquotesingle s invocation of Homology\textquotesingle s Mittag-Leffler lemma. It does not assume that inverse limits are exact in A. In particular it proves I is K-injective,\allowbreak{} but does not prove that an unrelated compatible map K\allowbreak{}→I is a quasi-isomorphism. The source explicitly keeps that latter issue open for a later lemma.

\subsection{3. The exact obstruction for a tower without split maps}\label{proofs-9746-10132-the-exact-obstruction-for-a-tower-without-split-maps}

DERIVED-UNDER-021 records a stronger conclusion of the preceding argument. Retain K-injective stages and existing termwise limits,\allowbreak{} but allow arbitrary transition maps. For each acyclic M use D\_\allowbreak{}n as above and put Z\_\allowbreak{}n\textasciicircum{}r\allowbreak{}=ker(d:\allowbreak{}D\_\allowbreak{}n\textasciicircum{}r\allowbreak{}→D\_\allowbreak{}n\textasciicircum{}\{r\allowbreak{}+1\})\allowbreak{}. For a tower A\_\allowbreak{}n of abelian groups write

Δ\_\allowbreak{}A:\allowbreak{}∏\_\allowbreak{}n A\_\allowbreak{}n\allowbreak{}→∏\_\allowbreak{}n A\_\allowbreak{}n,\allowbreak{} (x\_\allowbreak{}n)\allowbreak{}\allowbreak{}↦(x\_\allowbreak{}n-t\_\allowbreak{}n x\_\allowbreak{}\{n\allowbreak{}+1\})\allowbreak{},\allowbreak{} C\textasciicircum{}1(A\_\allowbreak{}n)\allowbreak{}\allowbreak{}=coker(Δ\_\allowbreak{}A)\allowbreak{}.

This defines the obstruction groups directly; no exactness of a limit functor is assumed. The exact sequences 0\allowbreak{}→Z\_\allowbreak{}n\textasciicircum{}\{-1\}\allowbreak{}→D\_\allowbreak{}n\textasciicircum{}\{-1\}\allowbreak{}→Z\_\allowbreak{}n\textasciicircum{}0\allowbreak{}→0 and exactness of products of abelian groups give,\allowbreak{} by the kernel-cokernel sequence for the vertical maps Δ,\allowbreak{} a natural identification

Hom\_\allowbreak{}K(M,\allowbreak{}I)\allowbreak{} ≅ ker(C\textasciicircum{}1(Z\_\allowbreak{}n\textasciicircum{}\{-1\})\allowbreak{}\allowbreak{}→C\textasciicircum{}1(D\_\allowbreak{}n\textasciicircum{}\{-1\})\allowbreak{})\allowbreak{}. (3.1)\allowbreak{}

Explicitly,\allowbreak{} for z\allowbreak{}=(z\_\allowbreak{}n)\allowbreak{}\allowbreak{}∈lim Z\_\allowbreak{}n\textasciicircum{}0 choose primitives h\_\allowbreak{}n with d h\_\allowbreak{}n\allowbreak{}=z\_\allowbreak{}n. Their discrepancies e\_\allowbreak{}n\allowbreak{}=h\_\allowbreak{}n-t\_\allowbreak{}n h\_\allowbreak{}\{n\allowbreak{}+1\} lie in Z\_\allowbreak{}n\textasciicircum{}\{-1\}. Changing the primitives by cycles changes (e\_\allowbreak{}n)\allowbreak{} by Δ\_\allowbreak{}Z of those cycles,\allowbreak{} and changing z by the boundary of a compatible h changes its class by zero. Moreover e\allowbreak{}=Δ\_\allowbreak{}D h,\allowbreak{} so its class maps to zero in C\textasciicircum{}1(D\_\allowbreak{}n\textasciicircum{}\{-1\})\allowbreak{}. Conversely,\allowbreak{} if a cycle family e represents an element of that kernel,\allowbreak{} write e\allowbreak{}=Δ\_\allowbreak{}D h. Then (d h\_\allowbreak{}n)\allowbreak{} is a compatible cycle family. If the class of e is zero in C\textasciicircum{}1(Z\_\allowbreak{}n\textasciicircum{}\{-1\})\allowbreak{},\allowbreak{} subtract the corresponding cycle family from h to make it compatible,\allowbreak{} so z is a boundary. These operations prove both the bijection and its naturality,\allowbreak{} including the exact map.

Thus I is K-injective precisely when the right side of (3.1)\allowbreak{} vanishes for every acyclic M. This is an explicit obstruction,\allowbreak{} not an assertion that arbitrary towers work. A sufficient condition is that for every acyclic M the tower D\_\allowbreak{}n\textasciicircum{}\{-2\} is Mittag-Leffler. Its quotient tower Z\_\allowbreak{}n\textasciicircum{}\{-1\}\allowbreak{}=im(D\_\allowbreak{}n\textasciicircum{}\{-2\}\allowbreak{}→D\_\allowbreak{}n\textasciicircum{}\{-1\})\allowbreak{} is then Mittag-Leffler too.

For completeness,\allowbreak{} a countable Mittag-Leffler tower A\_\allowbreak{}n has C\textasciicircum{}1(A\_\allowbreak{}n)\allowbreak{}\allowbreak{}=0. Let B\_\allowbreak{}n be the eventual image in A\_\allowbreak{}n; Mittag-Leffler makes it an attained image,\allowbreak{} and B\_\allowbreak{}\{n\allowbreak{}+1\}\allowbreak{}→B\_\allowbreak{}n is surjective by choosing a stage large enough for both images to have stabilized. The quotient tower A\_\allowbreak{}n/\allowbreak{}B\_\allowbreak{}n is pro-zero. On a pro-zero tower Q,\allowbreak{} the expression y\_\allowbreak{}n\allowbreak{}=Σ\_\allowbreak{}\{k≥n\}t\_\allowbreak{}\{n,\allowbreak{}k\}x\_\allowbreak{}k is a finite sum at each fixed n and satisfies Δ\_\allowbreak{}Q y\allowbreak{}=x. On a surjective tower B,\allowbreak{} solve Δ\_\allowbreak{}B y\allowbreak{}=x recursively,\allowbreak{} starting with any y\_\allowbreak{}1 and lifting y\_\allowbreak{}n-x\_\allowbreak{}n to y\_\allowbreak{}\{n\allowbreak{}+1\}. Given x in ∏A\_\allowbreak{}n,\allowbreak{} solve first in the quotient,\allowbreak{} lift a solution to ∏A\_\allowbreak{}n,\allowbreak{} then correct its remaining error in ∏B\_\allowbreak{}n. This proves Δ\_\allowbreak{}A is surjective. Applied to Z\_\allowbreak{}n\textasciicircum{}\{-1\},\allowbreak{} it makes the right side of (3.1)\allowbreak{} zero.

This sufficient condition includes some nonsplit towers. For K-injective I\_\allowbreak{}0,\allowbreak{}J let every stage be I\_\allowbreak{}0\allowbreak{}⊕J and every transition diag(id,\allowbreak{}0)\allowbreak{}. It is not termwise surjective where J is nonzero,\allowbreak{} but every Hom tower has stable image the I\_\allowbreak{}0 summand,\allowbreak{} and the termwise limit is I\_\allowbreak{}0. The actual projection and inclusion identify that limit and its K-injectivity. No novelty is claimed for this criterion or obstruction.

The source\textquotesingle s exact-adjoint preservation lemma is also valid. For v left adjoint to u with v exact,\allowbreak{} the degreewise adjunction carries chain maps and homotopies to one another,\allowbreak{} so Hom\_\allowbreak{}K(B)\allowbreak{}(M,\allowbreak{}uI)\allowbreak{}\allowbreak{}=Hom\_\allowbreak{}K(A)\allowbreak{}(vM,\allowbreak{}I)\allowbreak{}. An acyclic M remains acyclic under v:\allowbreak{} exactness preserves the kernel-image short exact sequences defining its cohomology. Therefore the last Hom group is zero. The existing article correction MC-STK-ERR-0884 is complete.

\subsection{4. Repairing the initial cutoff in the replacement lemma}\label{proofs-9746-10132-repairing-the-initial-cutoff-in-the-replacement-lemma}

The dimension function d is permitted to have value infinity. The limit condition n\allowbreak{}+d(K\textasciicircum{}n)\allowbreak{}\allowbreak{}→-infinity as n\allowbreak{}→-infinity implies finiteness far to the left,\allowbreak{} but permits finitely many infinite values in any remaining bounded range. Replacing only σ\_\allowbreak{}\{\textbackslash{}geq0\}K does not necessarily remove them. Choose instead an integer a such that d(K\textasciicircum{}n)\allowbreak{}<infinity for every n\textless a. Resolve σ\_\allowbreak{}\{\textbackslash{}geq a\}K by the source\textquotesingle s bounded below construction to a termwise monomorphic quasi-isomorphism into M,\allowbreak{} with M\textasciicircum{}n\allowbreak{}=0 for n\textless a and d(M\textasciicircum{}n)\allowbreak{}\allowbreak{}=0 for n≥a.

The class of d-zero objects contains zero. To check this from the stated axioms,\allowbreak{} choose A with d(A)\allowbreak{}\allowbreak{}=0 by axiom (1)\allowbreak{} and apply axiom (3)\allowbreak{} to 0\allowbreak{}→A\allowbreak{}→A\allowbreak{}→0\allowbreak{}→0; it gives d(0)\allowbreak{}≤max(-1,\allowbreak{}0)\allowbreak{}\allowbreak{}=0. Thus the cited bounded resolution lemma is applicable.

Splice the initial map into

…\allowbreak{}→K\textasciicircum{}\{a-2\}\allowbreak{}→K\textasciicircum{}\{a-1\}\allowbreak{}→M\textasciicircum{}a\allowbreak{}→M\textasciicircum{}\{a\allowbreak{}+1\}\allowbreak{}→…,\allowbreak{}

using the composite K\textasciicircum{}\{a-1\}\allowbreak{}→K\textasciicircum{}a\allowbreak{}→M\textasciicircum{}a at the join. Its square with the next differential is zero by the chain-map equation. The map from K is the identity below a and the resolution map at and above a. The short exact sequences with subcomplex σ\_\allowbreak{}\{\textbackslash{}geq a\} and common quotient σ\_\allowbreak{}\{\textbackslash{}leq a-1\} show it is a quasi-isomorphism. Equivalently the quotient of the monomorphic map is the acyclic quotient of that bounded resolution. All its d-values are now finite,\allowbreak{} d is zero for n≥a,\allowbreak{} and the far-left limit condition is unchanged.

DERIVED-NEW-033 changes the source\textquotesingle s fixed cutoff 0 to this cutoff a. Here is a complete example of the defect in the original finite-maximum argument. In abelian groups set d(G)\allowbreak{}\allowbreak{}=0 for divisible G and infinity otherwise. Finite sums and quotients of divisible groups are divisible. Consequently axioms (2)\allowbreak{} and (3)\allowbreak{} hold:\allowbreak{} the only nonautomatic case of (3)\allowbreak{} has both its first and middle groups divisible,\allowbreak{} hence its quotient is. Every group embeds into a divisible group. One direct construction adjoins roots to all its elements by a pushout of the monomorphism \allowbreak{}⊕\_\allowbreak{}\{(g,\allowbreak{}m)\allowbreak{}\}Z\allowbreak{}→\allowbreak{}⊕\_\allowbreak{}\{(g,\allowbreak{}m)\allowbreak{}\}Z whose components are multiplication by m≥1,\allowbreak{} with the map from the first sum to G sending its generator to g. Pushout preserves that monomorphism,\allowbreak{} so G embeds in the resulting group,\allowbreak{} where every original g has an m-th root. Iterating countably and taking the union embeds G in a divisible group. Thus axiom (1)\allowbreak{} holds as well. Let K be Z in degree -1 and zero elsewhere. It meets the far-left condition,\allowbreak{} but the source\textquotesingle s resolution of σ\_\allowbreak{}\{\textbackslash{}geq0\}K leaves it unchanged,\allowbreak{} so the printed maximum ξ is infinite. The corrected cutoff removes this term before the finite-maximum argument. The lemma remains true.

\subsection{5. Elementary replacements and a quantitative stabilization bound}\label{proofs-9746-10132-elementary-replacements-and-a-quantitative-stabilization-bound}

After Section 4 assume all d-values are finite and zero in sufficiently high degrees. For n with d(K\textasciicircum{}n)\allowbreak{}>0 choose a monomorphism j:\allowbreak{}K\textasciicircum{}n\allowbreak{}→M with d(M)\allowbreak{}\allowbreak{}=0. Use the pushout

M'\allowbreak{}=coker((j,\allowbreak{}-d\_\allowbreak{}K\textasciicircum{}n)\allowbreak{}:\allowbreak{}K\textasciicircum{}n\allowbreak{}→M\allowbreak{}⊕K\textasciicircum{}\{n\allowbreak{}+1\})\allowbreak{}.

The new differential into M is j d\_\allowbreak{}K\textasciicircum{}\{n-1\}; the next differential is m\allowbreak{}↦[m,\allowbreak{}0]\allowbreak{}; the differential out of M\textquotesingle{} is [m,\allowbreak{}x]\allowbreak{}\allowbreak{}↦d\_\allowbreak{}K\textasciicircum{}\{n\allowbreak{}+1\}x. The vertical map K\textasciicircum{}\{n\allowbreak{}+1\}\allowbreak{}→M' is x\allowbreak{}↦[0,\allowbreak{}x]\allowbreak{}. The identity {[}j y,\allowbreak{}0]\allowbreak{}\allowbreak{}=[0,\allowbreak{}d\_\allowbreak{}K\textasciicircum{}n y]\allowbreak{} proves the chain-map equation,\allowbreak{} and the relation is killed by the outgoing differential because d\_\allowbreak{}K²\allowbreak{}=0. All other components are unchanged.

Both changed vertical components are monomorphisms. The quotient complex is M/\allowbreak{}j(K\textasciicircum{}n)\allowbreak{} in degrees n and n\allowbreak{}+1,\allowbreak{} with identity differential:\allowbreak{} the second quotient M'/\allowbreak{}K\textasciicircum{}\{n\allowbreak{}+1\} is canonically M/\allowbreak{}j(K\textasciicircum{}n)\allowbreak{}. It is contractible. Hence this elementary replacement is a quasi-isomorphism. The defining short exact sequence for M\textquotesingle{} and the d axioms give exactly

d((K')\allowbreak{}\textasciicircum{}n)\allowbreak{}\allowbreak{}=0,\allowbreak{} d((K')\allowbreak{}\textasciicircum{}\{n\allowbreak{}+1\})\allowbreak{}≤max(d(K\textasciicircum{}n)\allowbreak{}-1,\allowbreak{}d(K\textasciicircum{}\{n\allowbreak{}+1\})\allowbreak{})\allowbreak{}.

If there are no positive d-values,\allowbreak{} the construction stops. Otherwise ξ\allowbreak{}=max\{n\allowbreak{}+d(K\textasciicircum{}n)\allowbreak{}:\allowbreak{}d(K\textasciicircum{}n)\allowbreak{}>0\} is a finite integer. Its maximum is attained at finitely many indices because weights tend to -infinity to the left and d is zero to the right. An elementary step cannot increase ξ. At a maximal-weight index n,\allowbreak{} any new positive weight at n\allowbreak{}+1 is at most ξ,\allowbreak{} and the old positive value at n disappears. With ξ fixed,\allowbreak{} the positive integer Σ\_\allowbreak{}\{n\allowbreak{}+d(K\textasciicircum{}n)\allowbreak{}\allowbreak{}=ξ\}(ξ-n)\allowbreak{} strictly decreases on every such step:\allowbreak{} a new maximal index can only replace n by n\allowbreak{}+1,\allowbreak{} or coincide with an already maximal index. Thus finitely many steps decrease ξ. Repeating either terminates or makes ξ tend to -infinity.

DERIVED-UNDER-022 gives a bound that also directly proves stabilization. Let d\_\allowbreak{}j be the finite dimensions immediately after the initial splice. For each integer m put

N\_\allowbreak{}m\allowbreak{}=Σ\_\allowbreak{}\{j:\allowbreak{}d\_\allowbreak{}j>0,\allowbreak{} j\allowbreak{}+d\_\allowbreak{}j≥m\}d\_\allowbreak{}j. (5.1)\allowbreak{}

This is a finite sum. During the algorithm let the same expression,\allowbreak{} formed from the current dimensions,\allowbreak{} be its current mass above level m. A replacement with n\allowbreak{}+d(K\textasciicircum{}n)\allowbreak{}≥m decreases that mass by at least one. Indeed it removes d(K\textasciicircum{}n)\allowbreak{}; if its recipient was already counted,\allowbreak{} the recipient increases by at most d(K\textasciicircum{}n)\allowbreak{}-1,\allowbreak{} and if it was not counted,\allowbreak{} its new counted value is at most d(K\textasciicircum{}n)\allowbreak{}-1. A replacement of lower weight cannot increase the mass:\allowbreak{} it cannot create a weight ≥m,\allowbreak{} and cannot increase a recipient already of weight ≥m.

A step changing a degree at least m has n\allowbreak{}+1≥m and d(K\textasciicircum{}n)\allowbreak{}≥1,\allowbreak{} hence n\allowbreak{}+d(K\textasciicircum{}n)\allowbreak{}≥m. Therefore at most N\_\allowbreak{}m elementary replacements affect the entire tail in degrees ≥m. This proves the source\textquotesingle s required local stabilization with an explicit bound.

Define L degree by degree by the eventual unchanged terms and differentials. Every component of K\allowbreak{}→L is the finite composite of the actual replacement maps affecting that component. Compatibility holds because the chain equations hold at every finite stage and eventually stabilize. Cohomology at degree r depends only on degrees r-1,\allowbreak{}r,\allowbreak{}r\allowbreak{}+1,\allowbreak{} which stabilize after finitely many steps; the finite composite to that stage is a quasi-isomorphism,\allowbreak{} so K\allowbreak{}→L induces an isomorphism on H\^{}r. Every surviving d-value is zero,\allowbreak{} since an eventual positive value would give a fixed weight while ξ tends to -infinity. No infinite colimit or exactness of an infinite limit is used. P02-E0076 also identifies the missing article at original 9968; its copyedit is separate from the cutoff repair.

\subsection{6. Finite cohomological dimension and the cone bound}\label{proofs-9746-10132-finite-cohomological-dimension-and-the-cone-bound}

Let F be left exact,\allowbreak{} with enough right acyclic objects,\allowbreak{} and suppose R\textasciicircum{}nF\allowbreak{}=0. If n\allowbreak{}=0,\allowbreak{} left exactness identifies R\^{}0F with F,\allowbreak{} so F is the zero functor. Its derived functor is zero on every complex,\allowbreak{} and all conclusions follow. In the remainder take n≥1.

The bounded below derived functor exists by enough acyclics. For a monomorphism A\allowbreak{}→A' with A\textquotesingle{} right acyclic and quotient C,\allowbreak{} the long exact sequence gives R\textasciicircum{}\{r\allowbreak{}+1\}F(A)\allowbreak{}≅R\textasciicircum{}rF(C)\allowbreak{} for r≥1. Starting with r\allowbreak{}=n and iterating shows R\textasciicircum{}rF\allowbreak{}=0 for all r≥n. Define

d(A)\allowbreak{}\allowbreak{}=max(\{0\}∪\{r:\allowbreak{}R\textasciicircum{}rF(A)\allowbreak{}≠0\})\allowbreak{}.

Then 0≤d(A)\allowbreak{}≤n-1. Additivity gives the direct-sum axiom for d. For 0\allowbreak{}→A\allowbreak{}→B\allowbreak{}→C\allowbreak{}→0 and r>max(d(A)\allowbreak{}-1,\allowbreak{}d(B)\allowbreak{})\allowbreak{},\allowbreak{} the two adjacent terms R\textasciicircum{}rF(B)\allowbreak{} and R\textasciicircum{}\{r\allowbreak{}+1\}F(A)\allowbreak{} in the long exact sequence vanish,\allowbreak{} so R\textasciicircum{}rF(C)\allowbreak{}\allowbreak{}=0. This proves the third d axiom. The enough-acyclics assumption supplies the first. The corrected replacement lemma therefore maps every complex to a complex of right acyclic objects. The admitted parentheses in MC-STK-ERR-0879 give this actual maximum.

Let L\allowbreak{}→M be a quasi-isomorphism between complexes of right acyclic objects. Fix a cohomological degree i and put c\allowbreak{}=i-n-1. Let Q\allowbreak{}=C(σ\_\allowbreak{}\{\textbackslash{}geq c\}L\allowbreak{}→σ\_\allowbreak{}\{\textbackslash{}geq c\}M)\allowbreak{}. The source uses stupid truncations. They agree with L,\allowbreak{}M in degrees at least c,\allowbreak{} so their cohomology agrees in degrees greater than c and vanishes below c. The long exact cone sequence therefore makes H\textasciicircum{}q(Q)\allowbreak{} zero except possibly for q\allowbreak{}=c-1,\allowbreak{}c.

Apply the bounded below hypercohomology spectral sequence to Q. Its second page has terms R\textasciicircum{}pF(H\textasciicircum{}q(Q)\allowbreak{})\allowbreak{},\allowbreak{} with 0≤p≤n-1 and q\allowbreak{}∈\{c-1,\allowbreak{}c\}. Every possible total degree p\allowbreak{}+q lies in

[c-1,\allowbreak{}c\allowbreak{}+n-1]\allowbreak{}\allowbreak{}=[i-n-2,\allowbreak{}i-2]\allowbreak{}. (6.1)\allowbreak{}

The filtration is bounded in each degree,\allowbreak{} so the same support holds for RF(Q)\allowbreak{}. In particular both H\textasciicircum{}\{i-1\}(RF(Q)\allowbreak{})\allowbreak{} and H\textasciicircum{}i(RF(Q)\allowbreak{})\allowbreak{} vanish. The triangle now proves an isomorphism on H\^{}i from RF(σ\_\allowbreak{}\{\textbackslash{}geq c\}L)\allowbreak{} to RF(σ\_\allowbreak{}\{\textbackslash{}geq c\}M)\allowbreak{}. Both are computed termwise by F because their terms are right acyclic and they are bounded below. Since i>c,\allowbreak{} their H\^{}i is that of F(L)\allowbreak{} and F(M)\allowbreak{}. Thus F(L)\allowbreak{}\allowbreak{}→F(M)\allowbreak{} is a quasi-isomorphism. The computation criterion constructs RF on all of D(A)\allowbreak{},\allowbreak{} and every complex of right acyclic objects computes it.

DERIVED-NEW-034 tightens the printed upper endpoint i-1 to i-2 in original 10048. The printed interval is a true but insufficient estimate:\allowbreak{} it does not exclude H\textasciicircum{}\{i-1\}(RF(Q)\allowbreak{})\allowbreak{},\allowbreak{} whose vanishing is needed for the claimed isomorphism. Formula (6.1)\allowbreak{} proves the missing vanishing rather than replacing it by an assumption.

\subsection{\texorpdfstring{7. Canonical truncation maps,\allowbreak{} their endpoints,\allowbreak{} and finite windows}{7. Canonical truncation  their  and finite windows}}\label{proofs-9746-10132-canonical-truncation-maps-their-endpoints-and-finite-windows}

For E\allowbreak{}∈D(A)\allowbreak{},\allowbreak{} the functor just constructed agrees with bounded below derivation on D\textasciicircum{}\allowbreak{}+(A)\allowbreak{}. In particular it sends D\^{}\{≥a\} into D\textasciicircum{}\{≥a\}:\allowbreak{} represent the input by a complex zero below a and use the bounded below acyclic resolution with the same lower bound. Apply RF to τ\_\allowbreak{}\{\textbackslash{}leq a\}E\allowbreak{}→E\allowbreak{}→τ\_\allowbreak{}\{\textbackslash{}geq a\allowbreak{}+1\}E. The long exact sequence proves H\textasciicircum{}i(RF(τ\_\allowbreak{}\{\textbackslash{}leq a\}E)\allowbreak{})\allowbreak{}\allowbreak{}→H\textasciicircum{}i(RF(E)\allowbreak{})\allowbreak{} is an isomorphism for i≤a. It is also a monomorphism for i\allowbreak{}=a\allowbreak{}+1. This proves the statement with the existing parentheses correction MC-STK-ERR-0878.

For the second source comparison write E as a complex L of right acyclic objects. For any integer c,\allowbreak{} RF(σ\_\allowbreak{}\{\textbackslash{}geq c\}L)\allowbreak{} is computed by σ\_\allowbreak{}\{\textbackslash{}geq c\}F(L)\allowbreak{}. The quantifier on c makes the existing unintroduced cutoff explicit; P02-E0079 is repaired this way,\allowbreak{} preserving the true general statement instead of replacing c throughout by one special value.

Set q\allowbreak{}=b-n. There are the two source sequences

H\textasciicircum{}q(L)\allowbreak{}[-q]\allowbreak{}\allowbreak{}→τ\_\allowbreak{}\{\textbackslash{}geq q\}L\allowbreak{}→τ\_\allowbreak{}\{\textbackslash{}geq q\allowbreak{}+1\}L,\allowbreak{} 0\allowbreak{}→im(d\_\allowbreak{}L\textasciicircum{}\{q-1\})\allowbreak{}[-q]\allowbreak{}\allowbreak{}→σ\_\allowbreak{}\{\textbackslash{}geq q\}L\allowbreak{}→τ\_\allowbreak{}\{\textbackslash{}geq q\}L\allowbreak{}→0.

For either left term T[-q]\allowbreak{},\allowbreak{} RF(T[-q]\allowbreak{})\allowbreak{} has cohomology only in degrees [q,\allowbreak{}q\allowbreak{}+n-1]\allowbreak{}\allowbreak{}=[b-n,\allowbreak{}b-1]\allowbreak{}. Their long exact sequences show that both maps RF(σ\_\allowbreak{}\{\textbackslash{}geq q\}L)\allowbreak{}\allowbreak{}→RF(τ\_\allowbreak{}\{\textbackslash{}geq q\}L)\allowbreak{} and RF(τ\_\allowbreak{}\{\textbackslash{}geq q\}L)\allowbreak{}\allowbreak{}→RF(τ\_\allowbreak{}\{\textbackslash{}geq q\allowbreak{}+1\}L)\allowbreak{} induce isomorphisms in degrees ≥b. The map RF(σ\_\allowbreak{}\{\textbackslash{}geq q\}L)\allowbreak{}\allowbreak{}→RF(L)\allowbreak{} is the actual inclusion σ\_\allowbreak{}\{\textbackslash{}geq q\}F(L)\allowbreak{}\allowbreak{}→F(L)\allowbreak{},\allowbreak{} also an isomorphism in degrees ≥b since n≥1. All these maps commute with L\allowbreak{}→τ\_\allowbreak{}\{\textbackslash{}geq q\allowbreak{}+1\}L. Thus that canonical map induces the asserted isomorphism

H\textasciicircum{}i(RF(E)\allowbreak{})\allowbreak{}\allowbreak{}→H\textasciicircum{}i(RF(τ\_\allowbreak{}\{\textbackslash{}geq b-n\allowbreak{}+1\}E)\allowbreak{})\allowbreak{},\allowbreak{} i≥b.

If the cohomology support of E lies in [a,\allowbreak{}b]\allowbreak{},\allowbreak{} the first comparison gives no RF cohomology below a and the second gives none at or above b\allowbreak{}+n. This proves support in [a,\allowbreak{}b\allowbreak{}+n-1]\allowbreak{},\allowbreak{} including unbounded endpoints.

DERIVED-UNDER-023 records the finite-window consequence. For every i,\allowbreak{} the canonical zigzag

E\allowbreak{}→τ\_\allowbreak{}\{\textbackslash{}geq i-n\allowbreak{}+1\}E←τ\_\allowbreak{}\{\textbackslash{}leq i\}τ\_\allowbreak{}\{\textbackslash{}geq i-n\allowbreak{}+1\}E

becomes isomorphisms on H\^{}i after applying RF. Hence H\textasciicircum{}i(RF(E)\allowbreak{})\allowbreak{} is canonically determined by the actual truncation window [i-n\allowbreak{}+1,\allowbreak{}i]\allowbreak{},\allowbreak{} retaining all its extension data,\allowbreak{} not merely the list of cohomology objects. The upper-truncation comparison is monomorphic at its next degree as already shown. The lower-truncation comparison E\allowbreak{}→τ\_\allowbreak{}\{\textbackslash{}geq c\}E is epimorphic on H\textasciicircum{}\{c\allowbreak{}+n-2\}(RF)\allowbreak{}:\allowbreak{} its triangle has first term RF(τ\_\allowbreak{}\{\textbackslash{}leq c-1\}E)\allowbreak{},\allowbreak{} supported at most in degree c\allowbreak{}+n-2,\allowbreak{} so the next cohomology of that term is zero. It is an isomorphism in degrees ≥c\allowbreak{}+n-1.

The left-derived assertions are obtained with the exact contravariant complex transformation D(K)\allowbreak{}\textasciicircum{}r\allowbreak{}=(K\textasciicircum{}\{-r\})\allowbreak{}\textasciicircum{}op and d\_\allowbreak{}D\textasciicircum{}r\allowbreak{}=(d\_\allowbreak{}K\textasciicircum{}\{-r-1\})\allowbreak{}\textasciicircum{}op used in PART14. It interchanges the canonical upper and lower truncations,\allowbreak{} sends H\^{}r to H\textasciicircum{}\{-r\},\allowbreak{} reverses arrows,\allowbreak{} and changes the interval [a,\allowbreak{}b]\allowbreak{} to [-b,\allowbreak{}-a]\allowbreak{}. The positive left-derived index n in this dual assertion means the vanishing in degree -n of LF(A[0]\allowbreak{})\allowbreak{},\allowbreak{} equivalently R\textasciicircum{}n(F\textasciicircum{}op)\allowbreak{}\allowbreak{}=0. The right interval [-b,\allowbreak{}-a\allowbreak{}+n-1]\allowbreak{} therefore returns exactly [a-n\allowbreak{}+1,\allowbreak{}b]\allowbreak{}. The canonical comparisons become

H\textasciicircum{}i(LF(τ\_\allowbreak{}\{\textbackslash{}leq a\allowbreak{}+n-1\}E)\allowbreak{})\allowbreak{}\allowbreak{}→H\textasciicircum{}i(LF(E)\allowbreak{})\allowbreak{},\allowbreak{} i≤a,\allowbreak{} H\textasciicircum{}i(LF(E)\allowbreak{})\allowbreak{}\allowbreak{}→H\textasciicircum{}i(LF(τ\_\allowbreak{}\{\textbackslash{}geq b\}E)\allowbreak{})\allowbreak{},\allowbreak{} i≥b.

These are the source\textquotesingle s maps with MC-STK-ERR-0880\textquotesingle s restored parenthesis. The dual finite window is [i,\allowbreak{}i\allowbreak{}+n-1]\allowbreak{}. At the first additional degree the upper-cut comparison is monomorphic; the lower-cut comparison is epimorphic at degree b-1. For n\allowbreak{}=0 all these functors are zero,\allowbreak{} as in the right-derived case. No replacement by a differently indexed derived functor has been made.

\subsection{\texorpdfstring{8. Receivers,\allowbreak{} dispositions,\allowbreak{} and the next section}{8.   and the next section}}\label{proofs-9746-10132-receivers-dispositions-and-the-next-section}

Current cohomology.tex 10300–\allowbreak{}10317 distinguishes K-injectivity of the termwise limit from its being a resolution. Sections 2–\allowbreak{}3 prove its K-injectivity invocation and preserve that distinction. Current injectives.tex 2078–\allowbreak{}2082 uses the existence of RF from enough K-injectives,\allowbreak{} exactly as Section 1 proves. Current perfect.tex 1643–\allowbreak{}1649 applies the exact-left-adjoint criterion,\allowbreak{} verified in Section 3. These are bounded checks of the stated invocations.

Current more-algebra.tex 23649–\allowbreak{}23678 applies the finite cohomological dimension theorem to inverse limits of abelian-group towers. The correction (6.1)\allowbreak{} preserves that theorem and its acyclic-complex computation. The exact obstruction groups in Section 3 were constructed directly and do not depend circularly on this receiving Rlim theorem. Current cohomology.tex 14218–\allowbreak{}14228 also uses that any complex of acyclics computes RF. That invocation is preserved. The preceding displayed identification of an ordinary Hom with RΓ at current 14213 is routed as a type-check observation for its own chapter review,\allowbreak{} not adjudicated here.

Ten more physical reports are complete:\allowbreak{} groups DERIVED-RECON-073 through 078. Eight are covered by existing units 0884,\allowbreak{} 0878,\allowbreak{} 0879,\allowbreak{} 0880; two need copyedit/\allowbreak{}notation operations (P02-E0076 and P02-E0079)\allowbreak{}. Totals are 135/\allowbreak{}176 reports,\allowbreak{} 78 groups,\allowbreak{} 124 already corrected,\allowbreak{} ten missing copyedit/\allowbreak{}notation reports and one missing mathematical-index report. There are 71 prior Derived units and 88 operations reconciled. NEW-033 and NEW-034 add two unreported findings; together with the two received copyedits there are four new operations,\allowbreak{} giving 65 cumulative operations. UNDER-021,\allowbreak{} 022 and 023 bring the editorial consequence total to 23.

Review is complete through original 10132; reading has reached 10310. Next is Derived colimits at 10139; unread source starts at 10311. The read-ahead functoriality and subsequence formulas require their own full diagram and report comparison before further operations are proposed. No new sessions,\allowbreak{} source composition,\allowbreak{} build,\allowbreak{} Lean run,\allowbreak{} cleanup,\allowbreak{} publication,\allowbreak{} or producer mathematical acceptance occurred.
